\section{Experiment description, analysis, and results}
\label{sec:results}

\subsection{Experimental Setup}

Two HVAC systems were selected for evaluation: System 1 (Air Handling Unit) and System 3
(Air Handling Unit). Both systems are drawn from real buildings with existing BACnet point
exports and engineering drawings available for ingestion. Both systems are installed on a real building site,
actual drawings of air handling units with real-world HVAC documentation, not synthetic test cases.

Four metrics are recorded for each AI agent run: total LLM token consumption (cumulative
context window size across all model calls in the session), estimated token cost in USD at
the applicable model pricing, wall-clock time from session start to validated TTL output,
and a qualitative assessment of the generated ASHRAE 223P graph based on completeness
(all equipment and sensors present), accuracy (correct topology and BACnet references),
and interpretability (graph readable by downstream applications).

The model used for both experiments is \textbf{Gemini 3.1 Pro Preview}, billed at
\$2.00/M tokens for prompts up to 200\,000 tokens and \$12.00/M for output and thinking
tokens (same prompt size tier).

\subsection{System 1: Air Handling Unit}

\subsubsection{Input Materials}

System 1 is a real-world Air Handling Unit (AHU) from an existing building installation.
The agent was provided three input artifacts: an HVAC engineering schematic drawing (PNG),
a BACnet point export CSV file containing raw control point data, and a control
documentation PDF. No additional context was given.

Figure~\ref{fig:system1-comparison} shows the original engineering drawing alongside the
AI-generated grid replication produced by the agent before ontology generation.

\begin{figure}[htbp]
  \centering
  \begin{minipage}[t]{0.48\textwidth}
    \centering
    \includegraphics[width=\textwidth]{system_one_original.png}
    \caption*{(a) Original engineering drawing}
  \end{minipage}
  \hfill
  \begin{minipage}[t]{0.48\textwidth}
    \centering
    \includegraphics[width=\textwidth]{system_one_ia_replicated.jpg}
    \caption*{(b) AI-generated grid replication}
  \end{minipage}
  \caption{System 1: original HVAC schematic (left) and the agent's replication on the
  Graphivac grid (right), produced as an intermediate step before ASHRAE 223P generation.}
  \label{fig:system1-comparison}
\end{figure}

\subsubsection{Performance Metrics}

\begin{table}[htbp]
  \caption{System 1 (AHU): AI Agent Performance}
  \label{tab:system1-ai}
  \centering
  \begin{tabular}{ll}
    \toprule
    \textbf{Metric} & \textbf{Value} \\
    \midrule
    Model used             & Gemini 3.1 Pro Preview \\
    Wall-clock time        & 17 min 11 s \\
    Peak context window    & 185{,}735 tokens \\
    Estimated token cost   & \$11.30 USD \\
    HVAC components placed & 21 (3 fans, 2 coils, 5 dampers, 6 sensors, \\
                           & \quad 1 filter, 1 humidifier, 2 VFDs, 1 valve) \\
    BACnet points mapped   & 44 points $\to$ 17 equipment instances \\
    Validation iterations  & 3 (auto-corrected by the agent) \\
    ASHRAE 223P graph      & 295 nodes, 949 relationships \\
    Session outcome        & Success \\
    \bottomrule
  \end{tabular}
\end{table}

\subsubsection{Graph Quality Assessment}

The agent successfully identified and placed all 21 HVAC components from the engineering
drawing, including fans, dampers, coils, and every instrumented sensor. The 44 BACnet
data points from the controls CSV were extracted and correctly associated with 17 equipment
instances, covering supply/return temperatures, damper positions, fan speeds, static
pressure, and humidity. Three minor code errors surfaced during ontology validation,
incorrect import names, a missing constructor parameter, and a Python operator syntax
issue, all of which the agent detected and corrected autonomously within the same
session, requiring no human intervention.

The resulting ASHRAE 223P graph contains 295 nodes and 949 relationships, encoding the
full equipment topology (inlet and outlet connection points, duct segments, medium types),
all BACnet external references, and the spatial composition of the AHU assembly.
Post-generation Cypher queries confirmed that fans, dampers, and sensor-to-BACnet
linkages were all correctly represented in the Neo4j database.

Figure~\ref{fig:system1-graph} shows the generated ASHRAE 223P semantic graph as
visualized in the SI-Mapper frontend.

\begin{figure}[htbp]
  \centering
  \begin{minipage}[t]{0.48\textwidth}
    \centering
    \includegraphics[width=\textwidth]{system_one_graph_full.jpg}
    \caption*{(a) Full graph view}
  \end{minipage}
  \hfill
  \begin{minipage}[t]{0.48\textwidth}
    \centering
    \includegraphics[width=\textwidth]{system_one_graph_detail.jpg}
    \caption*{(b) Detailed view of equipment topology}
  \end{minipage}
  \caption{System 1: generated ASHRAE 223P semantic graph in the SI-Mapper Neo4j
  visualization (295 nodes, 949 relationships). The detailed view shows equipment nodes
  with their connection points and BACnet external references.}
  \label{fig:system1-graph}
\end{figure}

\subsection{System 3: Air Handling Unit}

\subsubsection{Input Materials}

System 3 is a second real-world Air Handling Unit from a different building installation,
more complex than System 1. The agent received the same three input artifact types:
an HVAC engineering schematic (PNG), a BACnet point export CSV, and control documentation.

Figure~\ref{fig:system3-comparison} shows the original drawing and the AI-generated
grid replication.

\begin{figure}[htbp]
  \centering
  \begin{minipage}[t]{0.48\textwidth}
    \centering
    \includegraphics[width=\textwidth]{system_three.png}
    \caption*{(a) Original engineering drawing}
  \end{minipage}
  \hfill
  \begin{minipage}[t]{0.48\textwidth}
    \centering
    \includegraphics[width=\textwidth]{system_three_ia_replicated.jpg}
    \caption*{(b) AI-generated grid replication}
  \end{minipage}
  \caption{System 3: original HVAC schematic (left) and the agent's replication on the
  Graphivac grid (right).}
  \label{fig:system3-comparison}
\end{figure}

\subsubsection{Performance Metrics}

\begin{table}[htbp]
  \caption{System 3 (AHU): AI Agent Performance}
  \label{tab:system3-ai}
  \centering
  \begin{tabular}{ll}
    \toprule
    \textbf{Metric} & \textbf{Value} \\
    \midrule
    Model used             & Gemini 3.1 Pro Preview \\
    Wall-clock time        & 12 min 24 s \\
    Peak context window    & 121{,}962 tokens \\
    Estimated token cost   & \$10.85 USD \\
    HVAC components placed & 27 (3 fans, 3 VFDs, 3 heating coils, 1 cooling coil, \\
                           & \quad 2 filters, 3 dampers, 1 humidifier, 1 pump, \\
                           & \quad 1 valve, 9 sensors, 2 junctions) \\
    BACnet points mapped   & 34 points $\to$ 18 equipment instances \\
    Validation iterations  & 1 (auto-corrected by the agent) \\
    Session outcome        & Success \\
    \bottomrule
  \end{tabular}
\end{table}

\subsubsection{Graph Quality Assessment}

System 3 is a more complex AHU than System 1, featuring a heat recovery loop (3 heating
coils, a pump, and a 3-way valve) in addition to the standard cooling and ventilation
circuits. The agent correctly mapped all four air flow loops, supply, return, exhaust,
and heat recovery water, and placed all 27 components on the grid. The 34 BACnet
control points from the CSV were associated with 18 equipment instances, covering fan
speeds, coil temperatures, damper positions, filter differential pressure, and humidity.

Only one validation error arose during ontology generation: an API incompatibility in the
sensor-to-BACnet mapping syntax (the \texttt{@} operator was not supported directly on
pressure sensor objects). The agent identified the issue, applied the correct
\texttt{.observes @} pattern, and produced a clean TTL file without human intervention.

Figure~\ref{fig:system3-graph} shows the resulting ASHRAE 223P semantic graph in the
SI-Mapper frontend.

\begin{figure}[htbp]
  \centering
  \begin{minipage}[t]{0.48\textwidth}
    \centering
    \includegraphics[width=\textwidth]{system_three_graph_full.jpg}
    \caption*{(a) Full graph view}
  \end{minipage}
  \hfill
  \begin{minipage}[t]{0.48\textwidth}
    \centering
    \includegraphics[width=\textwidth]{system_three_graph_detail.jpg}
    \caption*{(b) Detailed view of equipment topology}
  \end{minipage}
  \caption{System 3: generated ASHRAE 223P semantic graph in the SI-Mapper Neo4j
  visualization. The detailed view shows equipment nodes with connection points, medium
  types, and BACnet external references.}
  \label{fig:system3-graph}
\end{figure}

\subsection{Human Engineer Baseline}

For comparison, an experienced HVAC engineer will perform the same task manually: reading
the same engineering drawings and BACnet point export CSV files, and producing an
equivalent ASHRAE 223P semantic model by hand using domain knowledge and available
ontology tooling. The human engineer is given the same inputs as the AI agent, no
additional context, to ensure a fair comparison of effort and output quality.

\begin{table}[htbp]
  \caption{Human Engineer Baseline}
  \label{tab:human-baseline}
  \centering
  \begin{tabular}{lll}
    \toprule
    \textbf{Metric} & \textbf{System 1} & \textbf{System 3} \\
    \midrule
    Time to complete
      & \todo{Fill after human trial}
      & \todo{Fill after human trial} \\
    Estimated labor cost (USD)
      & \todo{Fill after human trial}
      & \todo{Fill after human trial} \\
    ASHRAE 223P graph quality
      & \todo{Fill after human trial}
      & \todo{Fill after human trial} \\
    \bottomrule
  \end{tabular}
\end{table}

\subsection{Comparison: AI Agent vs.\ Human Engineer}

\begin{table}[htbp]
  \caption{AI Agent vs.\ Human Engineer: System 1 and System 3}
  \label{tab:comparison}
  \centering
  \begin{tabular}{llll}
    \toprule
    & \textbf{AI Agent} & \textbf{Human Engineer} & \textbf{Ratio} \\
    \midrule
    System 1: Time
      & 17 min 11 s
      & \todo{Human time}
      & \todo{Nx faster} \\
    System 1: Cost
      & \$11.30 USD
      & \todo{Human cost}
      & \todo{Nx cheaper} \\
    System 3: Time
      & 12 min 24 s
      & \todo{Human time}
      & \todo{Nx faster} \\
    System 3: Cost
      & \$10.85 USD
      & \todo{Human cost}
      & \todo{Nx cheaper} \\
    \bottomrule
  \end{tabular}
\end{table}

\todo{Write comparative analysis once both System 3 AI experiment and human baseline trials
are complete. Discuss time and cost ratios, graph quality comparison (completeness and
accuracy of BACnet mappings), cases where the AI struggled (e.g.\ ambiguous drawing
symbols), and implications for production deployment at scale.}
